\begin{center}\begin{tikzpicture}[font=\small,>=stealth,line join=round]
\draw[gray!40,rounded corners=3pt] (0,0) rectangle (14.2,8.3);\node at (7.1,7.700000000000001) {奇函数：把相加改成同一条曲线上的比较};
\begin{scope}[shift={(4.1,3.9499999999999997)},x=1.2cm,y=0.22999999999999998cm]\draw[->,gray] (-2.5,0)--(2.5,0) node[right]{$x$};\draw[->,gray] (0,-10)--(0,10) node[above]{$y$};\draw[blue,thick,solid,domain=-2.1:2.1,samples=180] plot(\x,{(\x)^3});\fill[orange!85!black] (1,1) circle(2.2pt);\fill[orange!85!black] (2,8) circle(2.2pt);\end{scope}\node[anchor=west,align=left,text width=5.0cm] at (8.4,4.0) {图取 $f(x)=x^3$\\[9pt]$a=2,\ b=-1$\\[9pt]$f(-b)=f(1)=1$\\[9pt]$f(a)+f(b)=8-1=7>0$\\[9pt]比较 $a=2>-b=1$};
\node[text width=12.6cm,align=center] at (7.1,.65) {图中标出的是比较用的 $a$ 和 $-b$；规律依赖奇性与严格单调性。};
\end{tikzpicture}\end{center}
